\section{A Scoped Simulation-Based Realization}
\label{sec:uc}
We now give an ideal functionality, an explicit simulator and a perfect
real/ideal argument for a restricted authorization-record service.
The result uses UC for static systems (UC4SS), with its fixed machine
interfaces and compatible composition operation \citep{canetti2025uc}.
It is a new written argument for the protocol defined here, not a claim
that the existing Python or Lean artifacts implement this functionality.
Its purpose is to establish a precise composable technical boundary;
the centralized simulation method is standard.

\paragraph{Scope of the realized service}
The construction has public request and response contents, fixed party
identities, a fixed static corruption set, one honest persistent coordinator
and session-local authenticated message channels. The coordinator implements
the application checks in ordinary local computation; there is no ideal
registry deciding whether a release is valid. Explicit privileged Tick
commands supply monotone logical time. The environment and adversary can
schedule messages and choose inputs adaptively, but cannot corrupt, reset
or roll back the coordinator. Root directory administration and the clock
endpoint are also outside the corruption set.

The application is a conservative ordinary-RELEASE specialization:
it retains attributed evidence, ordered approval, current authority,
exact record bindings, nonces, disclosure reservations and budgets.
It returns authorization and grant \emph{records}; it neither executes a
model nor issues an independently usable serving credential.
Mandatory institutional judgments remain attributed inputs.
The construction excludes adaptive corruption, confidential evidence,
global shared infrastructure and physical-time guarantees. A declaration
or manifest tag does not establish actual deployment identity.

\paragraph{Functionality and protocol}
Appendix~\ref{app:uc} defines the deterministic kernel $\Delta$, the
authenticated channel interface, the ideal functionality
$F_{\mathrm{MRA}}^\Delta$ and the real hybrid protocol
$\pi_\Delta^{\mathrm{auth}}$.
The functionality queues an attributed request and leaks its entire
contents. A simulator-controlled input-delivery callback executes the
kernel against \emph{then-current} state. The result is separately queued,
leaked and delivered by another callback. The real protocol consists of
client adapters, authenticated request/response channels and the honest
coordinator executing $\Delta$ atomically at request delivery.
These two delivery points are essential: a revocation may precede the
processing of an earlier-submitted grant request, or follow an already
committed grant whose receipt has not yet arrived.

\begin{uctheorem}[Perfect realization in the authenticated-channel hybrid]
\label{uc-thm:realization}
For the specified static interfaces, corruption set, polynomial resource
bounds and initialized kernel, for every admitted probabilistic
polynomial-time adversary $\mathcal A$ there is one polynomial-time
simulator $\mathcal S_{\mathcal A}$ such that, for every admitted
polynomial-time environment $\mathcal Z$ and auxiliary input $z$,
\[
\begin{split}
 &\Pr[\operatorname{Exec}(\pi_\Delta^{\mathrm{auth}},
                    \mathcal A,\mathcal Z,z)=1]\\
 &\quad=\Pr[\operatorname{Exec}(I_\Delta,
                    \mathcal S_{\mathcal A},\mathcal Z,z)=1].
\end{split}
\]
Here $I_\Delta$ is the ideal protocol for $F_{\mathrm{MRA}}^\Delta$,
including the dummy parties preserving the external interfaces.
Equality holds for every security parameter, not merely asymptotically.
\end{uctheorem}

The simulator runs $\mathcal A$, translates request and response handles,
and reproduces every allowed public notification. Its state invariant
equates the two application states, both pending queues and the honest
output histories. Induction over all interface actions proves identical
observable executions; the full proof is in Appendix~\ref{app:uc}.
The statistical failure budget $\alpha$ is absent from the emulation error:
both systems execute the same potentially fallible evidence procedures.
The separate scientific soundness premises are still required.

\begin{uctheorem}[Conditional transport substitution]
\label{uc-thm:transport}
Suppose a session-local transport protocol $\rho$ UC4SS-realizes the exact
authenticated-channel network used here, with error
$\epsilon_{\mathrm{tr}}$ and compatible interfaces, setup, leakage,
corruption and resource bounds. Then replacing that network in
$\pi_\Delta^{\mathrm{auth}}$ by $\rho$ realizes
$F_{\mathrm{MRA}}^\Delta$ with error at most $\epsilon_{\mathrm{tr}}$,
at the resources of the composed simulation.
\end{uctheorem}

This corollary is conditional on a transport realization, not a concrete
Ed25519 implementation theorem. Signature-based UC authentication is
established literature, with certification and session-management conditions
\citep{canetti2004authentication}. S1 alone does not discharge this
substitution premise. Globally verifiable signatures have additional
transferability issues \citep{canetti2016globalpki}; the result therefore
does not silently extend to government-wide shared keys or registries.
Neither theorem upgrades a recorded attestation into truth, nor proves
model confidentiality, a correct privacy observation boundary, effective
redress or faithful deployment.
